% ---------------------------------------------------------------
\section{Systematisation: Provenance Mechanisms as Primitives}\label{sec:sok}
% ---------------------------------------------------------------

\begin{figure}[t]
\centering
\resizebox{\columnwidth}{!}{% Fig. 1: pipeline stages (grey), mechanism primitives (blue, above), adversary capabilities (orange, below)
\begin{tikzpicture}[font=\sffamily\scriptsize, >={Stealth[length=1.7mm]},
  st/.style={draw=black!55, rounded corners=2pt, fill=black!5, minimum width=15mm,
             minimum height=8.5mm, align=center, inner sep=1pt, font=\sffamily\footnotesize},
  outb/.style={st, fill=white},
  mech/.style={draw=mechblue, fill=mechblue!9, text=mechblue!75!black, rounded corners=2.5pt, minimum height=3.6mm,
               inner sep=1.4pt, font=\sffamily\scriptsize\bfseries},
  adv/.style={draw=advorange, fill=advorange!9, text=advorange!80!black, rounded corners=1pt, minimum height=3.6mm,
              inner sep=1.4pt, font=\sffamily\scriptsize},
  fl/.style={->, black!70, line width=0.65pt}, tie/.style={black!35, line width=0.4pt}]
% stages
\node[st, minimum height=7mm] (src) at (0, 0.55) {Sources \ \textit{D}};
\node[st, minimum height=7mm] (gen) at (0,-0.55) {Generator \ \textit{S}};
\node[st] (tr)  at (2.35,0) {Transform\\\textit{F}};
\node[st] (cap) at (4.45,0) {Capture\\\textit{C}};
\node[st] (ex)  at (6.55,0) {Executor\\\textit{U}};
\node[outb, minimum width=12mm] (mdl) at (8.45,0) {Released\\model};
\node[outb, minimum width=22mm] (store) at (4.45,-2.15) {Record store\\\tiny (log or ledger)};
\node[outb, minimum width=15mm] (aud) at (7.5,-2.15) {Auditor};
% data flow
\draw[fl] (src.east) -- ++(0.3,0) |- (tr.west);
\draw[fl] (gen.east) -- ++(0.3,0) |- (tr.west);
\draw[fl] (tr) -- (cap); \draw[fl] (cap) -- (ex); \draw[fl] (ex) -- (mdl);
\draw[fl, dashed] (cap) -- node[right, font=\sffamily\tiny, text=black!60, pos=0.62]{record} (store);
\draw[fl, dashed] (store) -- (aud);
\draw[fl, dashed] (mdl.south) |- (aud.east);
% mechanisms
\node[mech] (m1) at (0,1.45) {WitSrc};          \draw[tie] (m1) -- (src);
\node[mech] (m2) at (2.35,1.05) {AttTrans};     \draw[tie] (m2) -- (tr);
\node[mech] (m3) at (4.45,1.05) {AttCap};       \draw[tie] (m3) -- (cap);
\node[mech, align=center] (m4) at (6.55,1.1) {BindFix /\\BindGrind}; \draw[tie] (m4) -- (ex);
\draw[mechblue, dashed, ->, line width=0.55pt] (m2.north) to[out=35,in=145] node[mech, solid, pos=0.5, fill=white] {WitIn} (m4.north);
\node[mech] (m5) at (2.35,-2.15) {Detector};    \draw[tie, ->] (m5) -- (store);
% adversary
\node[adv] (a1) at (0,-1.45) {AdvSource, AdvGen}; \draw[tie] (a1) -- (gen);
\node[adv] (a2) at (2.35,-1.0) {AdvTrans}; \draw[tie] (a2) -- (tr);
\node[adv] (a3) at (3.85,-1.0) {AdvCap};   \draw[tie] (a3) -- (cap);
\node[adv] (a4) at (6.55,-1.0) {AdvExec};  \draw[tie] (a4) -- (ex);
\node[adv] (a5) at (6.1,-2.95) {AdvWit, AdvTEE: trust base broken};
% legend
\node[mech, minimum width=3mm] (l1) at (-0.35,-2.95) {};
\node[right=0.4mm of l1, font=\sffamily\scriptsize] (t1) {mechanism};
\node[adv, minimum width=3mm, right=1.5mm of t1] (l2) {};
\node[right=0.4mm of l2, font=\sffamily\scriptsize] {adversary};
\end{tikzpicture}
}
\caption{The pipeline model. Grey: pipeline stages and the item sets of the model. Blue: where each mechanism primitive acts. Orange: what each adversary capability corrupts. Dashed: the record and the released model reach the auditor.}
\label{fig:pipeline}
\end{figure}

We group the published mechanisms by \emph{what they let an auditor check} rather than by
technology. Each group becomes one primitive of the model (Table~\ref{tab:sok}).

\paragraph{Computation binding}
Binding mechanisms let a verifier confirm that a model was produced by running a stated
training procedure on a stated dataset. Proof-of-learning~\cite{jia2021pol} has the prover
publish training checkpoints and the verifier re-execute selected segments. Its attacks
show that tolerance-based verification admits spoofed trajectories~\cite{fang2023pol}.
Zero-knowledge proofs of training give cryptographic soundness at substantial
cost~\cite{abbaszadeh2024zkpot,tan2025zkpot}. VeriLoRA extends them to LoRA fine-tuning of
LLMs with up to 13 billion parameters~\cite{liao2026verilora}. Exact re-execution is the
simplest binding. It is available whenever training is deterministic and affordable to
repeat, and we measure when it is (Section~\ref{sec:eval}). We separate binding that fixes
the training randomness in advance (\textsf{BindFix}) from binding that accepts the
prover's randomness (\textsf{BindGrind}), because the two differ in what they prove
(Theorems~\ref{thm:noinf} and~\ref{thm:grind}).

\paragraph{Source witnessing}
Signed manifests and receipts bind data to its supplier. VAMP signs dataset manifests with
the creator's key~\cite{stokes2021prevent}. Software supply-chain frameworks sign the steps
and artefacts of a build~\cite{torresarias2019intoto,newman2022sigstore}. The auditor's
check is that every recorded item is covered by a receipt (\textsf{WitSrc}). Receipts must
bind \emph{content}, not names: Carlini et al.\ show that web-scale datasets distributed
as URLs can be poisoned by changing what a URL serves between curation and
download (\emph{split-view} poisoning)~\cite{carlini2024poisoning}.

\paragraph{Input witnessing}
An observer at the training input, for example an attested data loader or the
external-input attestation of an enclave~\cite{huang2026exia}, reports what the executor
received. The auditor rejects a record that omits a reported input (\textsf{WitIn}).

\paragraph{Attested stages}
Trusted execution environments sign measurements that bind a stage's code, inputs and
outputs. Atlas runs attestation clients in TEEs and records every transformation in a
transparency log~\cite{spoczynski2025atlas}. Laminator issues hardware-attested property
cards for datasets, training and inference~\cite{duddu2025laminator}. We separate
attestation of the capture stage (\textsf{AttCap}) from attestation of transforms
(\textsf{AttTrans}), because they protect different relations.

\paragraph{Content detection}
Poisoning and backdoor detectors~\cite{goldblum2023dataset,tran2018spectral} inspect data
or trained models. PoisonSpot uses fine-grained training provenance to locate clean-label
backdoor samples~\cite{hailemariam2025poisonspot}. Data-use auditing tests whether given
data was used in training~\cite{huang2024datause,chen2025anonymity}. In the model a
detector is attested and examines the committed record (\textsf{Detector}). Its recall is
the free parameter $M$.

\paragraph{Record stores}
Certificate Transparency-style logs with co-signing witnesses give tamper evidence and
non-equivocation~\cite{rfc6962}. Hash chains give tamper evidence alone~\cite{li2025auditablellm}.
Blockchains add replicated, censorship-resistant inclusion~\cite{geren2025blockchain,mehammed2025fl}.
The model abstracts the store into properties $P_1$ (integrity), $P_2$ (non-equivocation)
and $P_8$ (inclusion liveness), and treats $P_1$ and $P_2$ as given
(Proposition~\ref{prop:ledger}).

\begin{table}[t]
\centering\scriptsize
\caption{Mechanism primitives, what the auditor checks, the trust they rest on, and
representative peer-reviewed work.}\label{tab:sok}
\begin{tabular}{@{}p{1.2cm}p{2.4cm}p{1.6cm}p{2.1cm}@{}}
\toprule
Primitive & Auditor's check & Rests on & Examples\\
\midrule
Bind & released model = Train(record) & re-execution or proof soundness & \cite{jia2021pol,abbaszadeh2024zkpot,tan2025zkpot,liao2026verilora}\\
WitSrc & every recorded item has a content receipt & sources' keys, non-collusion & \cite{stokes2021prevent,newman2022sigstore,torresarias2019intoto}\\
WitIn & record contains every observed training input & complete mediation, non-collusion & \cite{huang2026exia}\\
AttCap & record = capture input & TEE & \cite{spoczynski2025atlas,duddu2025laminator}\\
AttTrans & transform output $\subseteq$ inputs; synthetic report true & TEE & \cite{spoczynski2025atlas,duddu2025laminator}\\
Detector & no flagged poison in record & detector recall, TEE & \cite{tran2018spectral,hailemariam2025poisonspot,goldblum2023dataset}\\
Store & same unaltered record for all & witnesses or consensus & \cite{rfc6962,li2025auditablellm,mehammed2025fl}\\
\bottomrule
\end{tabular}
\end{table}

\paragraph{The base survey and what it leaves open}
Our starting point is the survey of Geren et al.~\cite{geren2025blockchain}, the most
complete peer-reviewed treatment of blockchain for LLM security and safety. It maps four
blockchain components (consensus, decentralised network, layer-2 zero-knowledge proofs and
smart contracts, and the verifiable Merkle-tree ledger~\cite{merkle1988signature}) onto a
taxonomy of LLM threats in training, prompting, data and supply chain. Its threat
categorisation follows Yao et al.~\cite{yao2024llmsecurity}. It builds on earlier
blockchain-for-AI reviews~\cite{salah2019blockchain,dinh2018ai} and on blockchain-enabled
federated learning~\cite{qu2023bcfl}. For the training phase, the survey collects systems that
use blockchain against poisoning and backdoors:
\begin{itemize}
\item a hash-chain-plus-proof-of-work framework presented as preventing data poisoning in
  smart power networks~\cite{keshk2020privacy};
\item a blockchain-based federated-learning framework that resists backdoors through robust
  aggregation~\cite{li2023dbfl};
\item DeepChain, which makes collaborative training auditable through a blockchain incentive
  mechanism~\cite{weng2021deepchain}.
\end{itemize}
The survey itself concludes that blockchain-based defences for LLMs still need comprehensive
integration and empirical validation. It also names data provenance as a direction in which
simple data-tracking solutions have remained absent, citing calls for provenance in
responsible AI~\cite{werder2022provenance}.

This paper takes up exactly those two open points. It does not catalogue further systems.
Instead, it asks what any combination of these mechanisms can \emph{prove}, and tests the
answer empirically. The model gives a precise reading of the survey's training-phase entries.
Hash chains and ledgers contribute record-store properties ($P_1$, $P_2$, $P_8$). They do not
prevent poison delivered by an authorised party (Theorem~\ref{thm:content}). Robust
aggregation is a content detector in our terms, so its guarantee is conditional on recall.
Auditability supports model-level claims only when the computation is bound to the record
(Theorem~\ref{thm:bind}). These readings rest on the survey's descriptions and the original
abstracts, not on re-implementations.

\begin{table*}[t]
\centering\scriptsize
\caption{Comparison with the closest prior work. ``Assurance claims'' and ``Adversary
assumptions'' follow each work's own description. ``Limitation'' is relative to the auditor's
question of this paper, not a judgement of the work's own aims.}\label{tab:compare}
\begin{tabular}{@{}p{2.1cm}p{2.2cm}p{2.4cm}p{2.3cm}p{3.1cm}p{3.4cm}@{}}
\toprule
Prior work & Mechanisms addressed & Assurance claims & Adversary assumptions & Limitation for the auditor's question & What this paper adds\\
\midrule
Geren et al.~\cite{geren2025blockchain} (survey) & Blockchain components (consensus, ledger, ZKP, smart contracts) & Taxonomy of how blockchain may help LLM security and safety & Threat categories, not a formal adversary & No sufficiency or minimality results; calls for empirical validation & Exact sufficiency and minimality per adversary; ledger's role reduced to $P_8$ (Prop.~\ref{prop:ledger}); empirical test\\
Lyu et al.~\cite{lyu2026survey}, Peng et al.~\cite{peng2026zkp} (surveys) & Verifiable / ZK ML techniques & Taxonomy of objectives, techniques and architectures; per-technique trust assumptions; Lyu: informal guidance on combining techniques & Per technique family, not a formal adversary & Combinations recommended informally, not decided against adversaries; Lyu leaves dataset completeness, label correctness and verifiable data poisoning open & Combinations decided exactly against six goals and 128 adversaries; $G_{\mathrm{compl}}$; content boundary (Thm.~\ref{thm:content})\\
PoL~\cite{jia2021pol}, attacks~\cite{fang2023pol} & Binding by spot-check replay & Model came from stated training, up to tolerance & Untrusted trainer & Tolerance admits spoofing; randomness chosen by prover & Grinding bound in terms of accepted outputs (Thm.~\ref{thm:grind}); measured on LLMs\\
zkPoT~\cite{abbaszadeh2024zkpot,tan2025zkpot}, VeriLoRA~\cite{liao2026verilora} & Cryptographic binding & Sound proof of training & Untrusted trainer & Lineage and content outside scope; Kaizen and VeriLoRA accept prover-chosen initialisation and randomness; Tan et al.\ show seed rejection sampling and remove randomness by optimum vicinity, for convex objectives only & Shows what binding needs beside it (Thms.~\ref{thm:auth},~\ref{thm:content}); BindFix vs BindGrind; influence bound for non-convex training (Thm.~\ref{thm:grind})\\
Atlas~\cite{spoczynski2025atlas}, Laminator~\cite{duddu2025laminator} & TEE attestation, transparency logs & Tamper detection, attested lineage and properties & MLaaS providers untrusted; TEEs and verifiers trusted & Training on GPUs outside TEE (Atlas); content out of scope & Which goals survive without attested training or under TEE compromise\\
VAMP~\cite{stokes2021prevent}, Sigstore~\cite{newman2022sigstore}, in-toto~\cite{torresarias2019intoto} & Signed manifests and supply-chain attestations & Authenticity and integrity of artefacts; VAMP: poisoning prevention & Infrastructure trusted; signers authorised & Authorised poison passes & Formal and empirical counterexample (Thm.~\ref{thm:content}, G3)\\
CT logs~\cite{rfc6962}, ledgers~\cite{mehammed2025fl,li2025auditablellm} & Record stores & Integrity, non-equivocation, inclusion & Log operator untrusted, witnesses honest & Say nothing about what enters the record & Store properties separated from pipeline goals\\
Detectors~\cite{tran2018spectral,hailemariam2025poisonspot}, data-use audit~\cite{huang2024datause} & Content inspection & Detect poisoned or used items & Varies & No transfer to the released model without binding & Conditional content claim (Thm.~\ref{thm:content}); recall and precision measured\\
\bottomrule
\end{tabular}
\end{table*}

\paragraph{Worked example}
VAMP~\cite{stokes2021prevent} binds datasets, training software and models to signed
manifests, and its title presents this as \emph{preventing} poisoning. In our model its
deployment corresponds to \textsf{WitSrc} under a trusted training infrastructure: the
trans-only profile, filter-only pipeline. Table~\ref{tab:minimal} then gives lineage
($G_{\mathrm{auth}}$), and the scheme is sound for that claim. The same table shows that no
mechanism set achieves $G_{\mathrm{zero}}$. The proof of Theorem~\ref{thm:content} gives the
counterexample in one line: a dataset creator holding a valid signing key signs a poisoned
item. The taxonomies we reviewed record that VAMP uses signatures and provenance; they do not
make this boundary explicit. Our experiment G3 (Section~\ref{sec:eval}) realises it on
DistilBERT. Every poisoned example carried a valid receipt, and the backdoor raised attack
success by 20 to 75 points.

\paragraph{Other surveys}
Lyu et al.~\cite{lyu2026survey} and Peng et al.~\cite{peng2026zkp} survey verifiable ML and
zero-knowledge ML. Together with the base survey, these works classify techniques and their
trust assumptions. Lyu et al.\ go further than a list of techniques: they tabulate the trust
assumptions of each technique family and give practitioners informal guidance on which
techniques to combine for a use case. They also state three gaps that bear on our question.
No surveyed system gives model binding, update integrity and supply-chain provenance in one
proof, and the cost of combining them is open. Label correctness, sampling fidelity and
dataset completeness lie outside the data-integrity guarantees they survey. Verifiable data
poisoning is listed as an open challenge. Our model addresses these gaps directly: combinations
are decided against a formal adversary (Table~\ref{tab:minimal}), completeness is a goal of its
own ($G_{\mathrm{compl}}$), and Theorem~\ref{thm:content} states what can and cannot be
verified about poisoning. Our contribution is complementary: a model in which combinations of those
techniques are compared against stated adversaries and goals, with exact sufficiency and
minimality results, and an experimental test of the model. The structured search behind this
section is described in Section~\ref{sec:lit}.
